\ssscp{Distribution of words for ‘laugh’}{15-04-04-01}
From the lexicon,
the strongest etymon in support of CEMP is **malip ‘laugh, smile’.
\frf{15-08} shows its distribution in the region,
along with known pockets of languages that get their word for
‘laugh’ from other sources,
including PMP *tawa ‘laugh’.
White circles indicate languages investigated reflecting etyma
other than **malip or PMP *tawa for ‘laugh’.

As seen in \frf{15-08},
reflexes of PMP *tawa ‘laugh’ persist across \tsc{Bima-Lembata}
languages, with only isolated reflexes of **malip ‘laugh,
smile’ in that subgroup.
\tsc{Taliabu} languages appear to get their words for ‘laugh’
from neither *tawa nor **malip.
While the distribution of reflexes of **malip east of the Blust
line extending into Oceania is quite clear
and remarkable, it is nevertheless not a replacement innovation
for *tawa in all Wallacean subgroups,
but it is a replacement innovation in many.

\begin{figure}[h]\centering
	\includegraphics[width=\columnwidth]{laugh.pdf}
	\fcp{Distribution of words for ‘laugh’}{15-08}
\end{figure}


